\chapter{Đồng Tiền Có Mồ Hôi}
\markboth{Đồng Tiền Có Mồ Hôi}{Every Coin Carries Sweat}

\section{Khi Một Món Đồ Đẹp Đi Qua Trước Mắt}
\begin{truyen}{Khi Một Món Đồ Đẹp Đi Qua Trước Mắt}{When Something Attractive Appears Before Your Eyes}
\chuhoa{C}{uối một ngày dài, người thầy trở về nhà với lưng áo còn dính mồ hôi. Vừa ngồi xuống chưa lâu, điện thoại hiện lên một món đồ giảm giá rất hấp dẫn.}
\textit{(At the end of a long day, the teacher returned home with sweat still on his back. He had barely sat down when his phone displayed an attractive discounted item.)}

Không phải anh chưa từng thiếu nó. Cũng không phải cuộc sống sẽ tốt hơn nhiều nếu mua ngay. Nhưng sau những giờ đứng lớp, những buổi dạy thêm, những đoạn đường đi về trong mệt mỏi, người ta thường dễ mềm lòng với bản thân và gọi sự chiều mình là phần thưởng.
\textit{(It was not that he had been missing it terribly, nor that life would improve greatly if he bought it immediately. But after hours of teaching, tutoring, and traveling home in exhaustion, people often soften toward themselves and call self-indulgence a reward.)}

Anh cầm điện thoại rất lâu. Một bên là sự thích thú nhất thời. Một bên là cảm giác nhớ đến những giờ lao động âm thầm không ai thấy. Cuối cùng, anh tắt màn hình, rót một cốc nước lọc, rồi nói nhỏ với chính mình rằng thứ cần được thưởng không phải lòng ham, mà là sự bình tĩnh.
\textit{(He held the phone for a long time. On one side stood brief excitement. On the other stood the memory of silent working hours no one sees. In the end, he turned off the screen, poured a glass of water, and quietly told himself that what deserved reward was not desire, but calmness.)}
\end{truyen}

\begin{baihoc}
Người biết quý công sức lao động sẽ không tiêu tiền như thể tiền tự nhiên mà có.
\textit{(A person who values labor will not spend money as though it appeared by itself.)}
\end{baihoc}

\begin{ghinhoanh}
Money is slow labor made visible.

Respect the work behind the coin.
\end{ghinhoanh}

\begin{tuhocnhanh}
1. Em thường yếu lòng nhất vào lúc nào? \textit{When are you usually most vulnerable to spending?}

2. Em có đang gọi ham muốn là phần thưởng không? \textit{Are you calling desire a reward?}
\end{tuhocnhanh}
\ngancach

\section{Điển Tích Gia Cát Lượng}
\begin{truyen}{Kiệm Dĩ Dưỡng Đức}{Frugality Nurtures Virtue}
\chuhoa{G}{ia Cát Lượng để lại một câu rất ngắn nhưng rất sâu: ``Tĩnh dĩ tu thân, kiệm dĩ dưỡng đức.'' Người xưa không xem tiết kiệm là mẹo giữ tiền, mà là nếp giữ mình.}
\textit{(Zhuge Liang left a brief but profound line: ``Cultivate the self through calmness; nurture virtue through frugality.'' The ancients did not regard thrift merely as a trick for keeping money, but as a way of keeping oneself.)}

Tiết kiệm giúp con người chậm lại trước cám dỗ. Khi chậm lại, ta bắt đầu nhìn rõ hơn cái gì là nhu cầu, cái gì chỉ là nôn nóng, sĩ diện, buồn chán hay muốn tự an ủi. Một người giữ được cái chậm ấy thường cũng giữ được nhiều thứ khác: lời nói, lòng giận, và cả những quyết định quan trọng.
\textit{(Frugality helps a person slow down before temptation. In that slowness, we begin to see more clearly what is a need and what is merely impatience, pride, boredom, or self-comfort. Whoever preserves that slowness often preserves other things too: speech, anger, and important decisions.)}

Thế nên giữ tiền không chỉ là chuyện của chiếc ví. Nó là chuyện của nhân cách. Người phung phí lâu ngày không chỉ mất tiền. Họ còn dễ đánh mất cảm giác biết đủ, biết sợ và biết ơn.
\textit{(So keeping money is not only a matter of the wallet. It is a matter of character. Those who waste for a long time lose more than money. They also risk losing the sense of enoughness, caution, and gratitude.)}
\end{truyen}

\begin{baihoc}
Giữ đồng tiền cho đúng là một cách giữ nếp sống cho ngay.
\textit{(Keeping money wisely is one way of keeping one’s life in order.)}
\end{baihoc}

\begin{ghinhoanh}
Frugality guards more than money.

It also guards character.
\end{ghinhoanh}

\begin{tuhocnhanh}
1. Khi tiêu tiền, em đang nghe lý trí hay nghe cảm xúc? \textit{When spending, are you listening to reason or emotion?}

2. Thói quen nào đang làm em mềm lòng nhất? \textit{Which habit makes you weakest?}
\end{tuhocnhanh}
\ngancach

\section{Tự Răn Mình}
\begin{truyen}{Một Câu Tự Nhắc Mỗi Tối}{One Nightly Sentence of Self-Reminder}
\chuhoa{T}{ừ hôm ấy, người thầy đặt ra cho mình một câu ngắn: ``Đồng tiền này đã đổi bằng bao nhiêu giờ đời mình?'' Mỗi lần chuẩn bị mua một món chưa thật sự cần, anh đều tự hỏi như vậy.}
\textit{(From that day on, the teacher kept one short sentence for himself: ``How many hours of my life paid for this coin?'' Every time he was about to buy something unnecessary, he asked himself that question.)}

Ban đầu câu hỏi ấy làm anh khó chịu, vì nó cắt ngang sự hứng khởi. Nhưng dần dần, chính sự khó chịu ấy cứu anh khỏi nhiều lần tiêu theo cảm xúc. Câu tự răn không làm đời nhẹ ngay. Nó chỉ làm bàn tay bớt vội hơn. Và nhiều khi, bớt vội là đã bớt sai rất nhiều.
\textit{(At first the question annoyed him because it interrupted excitement. But slowly, that discomfort saved him from many emotional purchases. A self-reminder does not make life easier at once. It simply makes the hand less hasty. And often, less haste already means far fewer mistakes.)}
\end{truyen}

\begin{baihoc}
Người biết dừng một nhịp trước khi tiêu tiền thường giữ được nhiều tháng bình yên sau đó.
\textit{(A person who pauses for one breath before spending often keeps many peaceful months afterward.)}
\end{baihoc}

\begin{ghinhoanh}
Pause before spending.

One breath can protect many days.
\end{ghinhoanh}

\begin{tuhocnhanh}
1. Câu tự nhắc nào em muốn viết cho mình? \textit{What self-reminder do you want to write for yourself?}

2. Nếu dừng lại một nhịp, món đồ em đang muốn mua còn quan trọng như cũ không? \textit{If you pause for a moment, does the item still matter as much?}
\end{tuhocnhanh}

\begin{turanminh}
1. Tôi đã từng lớn lên trong thiếu thốn, nên không được phép xem nhẹ đồng tiền do mồ hôi mình làm ra.

2. Mỗi lần muốn tự thưởng bằng một món chưa cần, tôi phải nhớ đến những buổi đứng lớp, những giờ dạy thêm và những phần đời mình đã đổi lấy khoản tiền ấy.

3. Tôi không tiết kiệm vì sợ sống khổ. Tôi tiết kiệm để con mình đỡ phải khổ như mình đã từng.
\end{turanminh}

\begin{dayhoc}
1. Có thể hỏi học sinh: ``Nếu món đồ em muốn mua phải đổi bằng công sức lao động của cha mẹ trong một ngày, em còn muốn nó ngay lúc này không?''

2. Nhắc các em rằng biết quý tiền bắt đầu từ biết quý công sức, quý sách vở, quý bữa cơm và quý thời gian của người nuôi mình lớn.

3. Khi dạy trên lớp, thầy có thể chốt bằng một câu ngắn: ``Đồng tiền nào cũng có mồ hôi ở sau.''
\end{dayhoc}